% Capítulo 1: Introducción
% Propósito: Describir los antecedentes, planteamiento del problema, justificación, formulación y los objetivos (general y específicos) del proyecto de grado.


\section{Introducción}

El transporte terrestre de carga forma parte del movimiento diario de mercancías entre empresas, ciudades y diferentes puntos de distribución. Con el paso del tiempo, la forma de administrar estas operaciones también ha cambiado, especialmente por la incorporación de tecnologías que permiten manejar una mayor cantidad de información durante los procesos logísticos. En este sentido, la logística inteligente ha comenzado a apoyarse en herramientas digitales para automatizar actividades, monitorear envíos y mejorar la organización de las operaciones \parencite{kalkha2023rising}.

Este avance tecnológico ha hecho que la gestión del transporte ya no se observe únicamente desde el traslado físico de una mercancía. Actualmente también se estudian aspectos relacionados con el seguimiento de los envíos y la planificación de los recorridos. \textcite{zhang2024ecommerce}, por ejemplo, desarrollaron una propuesta en la que el seguimiento de paquetes y la planificación de rutas forman parte de una misma solución logística. Por otro lado, \textcite{xin2022logistics} muestran que la búsqueda de recorridos más adecuados continúa siendo un tema de interés dentro de la distribución de mercancías, debido a su relación con el uso de los vehículos, los tiempos de recorrido y la organización de las entregas.

Tomando como referencia esta evolución de los sistemas logísticos surge TravelDLS, un proyecto orientado al transporte terrestre de carga pesada. La propuesta consiste en desarrollar una plataforma web que permita reunir diferentes empresas transportistas dentro de un mismo entorno y facilitar la relación entre empresas, conductores y clientes. Además de servir como espacio para administrar las operaciones propias de cada empresa, se busca que las empresas transportistas puedan dar a conocer sus servicios y tener mayor visibilidad ante personas o negocios que necesiten contratar transporte de carga. El sistema también incorporará funciones relacionadas con los viajes, la ubicación de las unidades y la planificación de recorridos, manteniendo cada una de estas funciones conectada con el servicio que se está realizando.

A partir de esta propuesta, el presente proyecto se encuentra organizado en diferentes apartados que permiten desarrollar de manera progresiva la propuesta de TravelDLS. En primer lugar, se presenta el \textbf{\hyperref[sec:planteamiento]{Planteamiento del Proyecto}}, donde se describe la situación que da origen a la propuesta y las necesidades que se pretenden atender. Posteriormente, se establecen los \textbf{\hyperref[sec:objetivos]{Objetivos}} que orientan el desarrollo del proyecto y se presentan los \textbf{\hyperref[sec:antecedentes]{Antecedentes}}, los cuales permiten conocer investigaciones y soluciones relacionadas con la gestión del transporte, el seguimiento y la optimización de rutas. A continuación, la \textbf{\hyperref[sec:justificacion]{Justificación}} expone la importancia de la propuesta, mientras que el \textbf{\hyperref[sec:alcance]{Alcance}} y las \textbf{\hyperref[sec:limitaciones]{Limitaciones}} delimitan las funciones que serán desarrolladas y los aspectos que quedan fuera de esta etapa. Seguidamente, se analiza la \textbf{\hyperref[sec:viabilidad]{Viabilidad del Proyecto}} desde las perspectivas técnica, operativa, económica y legal. El \textbf{\hyperref[sec:marco-teorico]{Marco Teórico}} reúne los antecedentes del proyecto, los conceptos y las bases teóricas necesarias para comprender los componentes de TravelDLS. Finalmente, la \textbf{\hyperref[sec:metodologia]{Metodología}} describe el procedimiento seguido para la selección de investigaciones, la recolección y análisis de información y la organización del proceso de desarrollo de la plataforma.

\subsection{Planteamiento del Proyecto}
\label{sec:planteamiento}

El transporte terrestre de carga pesada implica mucho más que mover una mercancía de un punto a otro. En cada servicio también deben coordinarse vehículos, conductores, cargas, recorridos, tiempos y comunicación con el cliente. Cuando esa información se encuentra separada o no se actualiza de forma oportuna, resulta más difícil saber qué está ocurriendo durante el viaje y reaccionar ante cambios, atrasos o inconvenientes. Esta necesidad de mantener una mejor comunicación también aparece en estudios sobre logística, donde se señalan dificultades entre clientes, conductores y empresas, además de la necesidad de conocer el estado de los envíos durante el traslado \parencite[p. 2]{sultan2022real}.

El seguimiento y la planificación de rutas son dos apoyos que pueden ayudar a ordenar este tipo de operaciones. \textcite{pauzuoliene2024smart} encontraron que, en una muestra de 46 empresas logísticas de Lituania, el 69.6 \% utilizaba tecnologías para dar seguimiento a la carga en tiempo real y el 60.9 \% empleaba programas para planificar o seleccionar rutas (p. 988). En Centroamérica también se han señalado necesidades de digitalización y coordinación. En Panamá, por ejemplo, se han identificado barreras tecnológicas, baja digitalización y dificultades de coordinación dentro del transporte terrestre \parencite[p. 10]{castillo2025potenciando}. Aunque estos estudios corresponden a otros contextos, ayudan a entender que mantener conectada la información del viaje sigue siendo una necesidad presente en distintas operaciones de transporte.

En Nicaragua ya existen trabajos que atienden partes de esta situación. \textcite{aragon2018implementacion} desarrollaron un sistema para administrar fletes terrestres con la intención de ordenar los controles y facilitar el acceso a la información logística (p. 2). Por su parte, \textcite{garcia2022desarrollo} abordaron en León la necesidad de conocer recorridos y ubicación de unidades de transporte mediante una aplicación con funciones de geolocalización en tiempo real (pp. 3, 7). Son proyectos diferentes a TravelDLS, pero muestran que tanto la organización de la información como la ubicación de los vehículos ya han sido necesidades abordadas desde la informática en el país.

A partir de este panorama, TravelDLS se concentra en una necesidad más específica: que las empresas dedicadas al transporte terrestre de carga pesada puedan organizar sus operaciones y, al mismo tiempo, mantener conectada la información que necesitan los conductores y los clientes durante el servicio. El problema no consiste únicamente en registrar un viaje. También implica relacionar el vehículo, el conductor, la carga, la ruta, la ubicación y el estado del servicio para que la información tenga sentido durante todo el recorrido.

Por ello, la pregunta que guía el proyecto es la siguiente: ¿Cómo mejorar la gestión, coordinación, localización, seguimiento y planificación de rutas de las operaciones de transporte terrestre de carga pesada, de manera que empresas transportistas, conductores y clientes dispongan de información centralizada y oportuna durante el servicio?

Como respuesta a esta necesidad se propone TravelDLS, una plataforma web bajo el modelo de Software como Servicio (SaaS). La idea es reunir en un mismo entorno la información de empresas transportistas, vehículos, conductores, cargas, viajes, rutas y clientes, de manera que las distintas partes del servicio no queden aisladas. La plataforma también contempla la publicación y contratación de servicios, el seguimiento de los viajes, la consulta del estado de la carga y un módulo para apoyar la planificación de rutas con varios destinos. TravelDLS no pretende reemplazar las decisiones de la empresa, sino ofrecer un medio más ordenado para consultar información y apoyar la gestión diaria de los viajes.

\subsection{Objetivos}
\label{sec:objetivos}

\subsubsection{Objetivo General}
 Desarrollar un sistema web denominado TravelDLS para la gestión, localización y trazabilidad en tiempo real de las operaciones de transporte terrestre de carga pesada, incorporando la optimización matemática de rutas para fortalecer los procesos administrativos y logísticos de las empresas transportistas y brindar mayor transparencia a los clientes durante el traslado de sus cargas.

\subsubsection{Objetivos Específicos}
\label{sec:objetivos-especificos}
\begin{itemize}
	\item Analizar los requerimientos funcionales y no funcionales del sistema, evaluando y comparando métodos de optimización de rutas basados en el Problema del Agente Viajero (Vecino Más Cercano, 2-opt y Google OR-Tools) para seleccionar la solución más eficiente.
	\item Diseñar la arquitectura del sistema, incluyendo la base de datos y la estructura de los módulos que conformarán la plataforma.
	\item Proponer un prototipo de plataforma web (backend y frontend), implementando el módulo de optimización con Google OR-Tools para la gestión eficiente de usuarios, vehículos, conductores, cargas, viajes y rutas.
\end{itemize}

\subsection{Antecedentes}
\label{sec:antecedentes}

Para entender mejor de dónde parte TravelDLS, se revisaron investigaciones realizadas de manera Internacional como Nacional que abordan problemas parecidos desde distintos puntos de vista. Algunos trabajos se concentran en el seguimiento de la carga, otros en la planificación de rutas y otros en la administración de los servicios de transporte. Revisarlos permite reconocer qué soluciones se han probado, qué resultados han obtenido y qué ideas pueden servir como referencia para el proyecto, sin asumir que una propuesta aplicada en otro contexto funcionará de la misma manera en TravelDLS.

\phantomsection\label{ant:pauzuoliene}
En el ámbito internacional, \textcite{pauzuoliene2024smart} realizaron la investigación titulada \textbf{Smart technologies integration and challenges in the context of logistics companies}, en donde estudiaron el uso de tecnologías inteligentes en empresas logísticas de Lituania. La investigación contó con 46 representantes del sector y mostró que el 69.6 \% utilizaba seguimiento de carga en tiempo real y el 60.9 \% empleaba programas para planificar o seleccionar rutas. Los autores también relacionan estas herramientas con la posibilidad de conocer la ubicación, el estado de la carga y los tiempos esperados de entrega (pp. 988-989). Para TravelDLS, este trabajo es útil porque muestra que el seguimiento y la planificación pueden formar parte de una misma gestión logística, aunque el estudio se enfoca en la adopción de tecnologías y no en el desarrollo de una plataforma para carga pesada.

\phantomsection\label{ant:sultan}
También a nivel internacional, \textcite{sultan2022real} llevaron a cabo el estudio titulado \textbf{A real-time mobile and web-based logistics management system for collaborative urban freight transportation} en el cual propusieron un sistema móvil y web para apoyar el transporte urbano de mercancías. La solución reúne funciones de registro, reservas, conductores, administración, inventario, seguimiento de pedidos y planificación de rutas (pp. 1-3). Los autores además señalan problemas de comunicación entre clientes, conductores y empresas, junto con la necesidad de que el cliente pueda conocer el estado de su envío durante el traslado (p. 2). Este trabajo se acerca a TravelDLS por la forma en que integra varios participantes y procesos dentro de una misma solución, aunque TravelDLS estará dirigido específicamente al transporte terrestre de carga pesada y a varias empresas usuarias.

\phantomsection\label{ant:aguilar}
Al acercar la revisión a Centroamérica, \textcite{aguilar2022optimizacion}, en Honduras,publicaron el artículo titulado \textbf{Optimización de Rutas de compañías distribuidoras y surtidoras en Tegucigalpa} en donde se estudió la optimización de rutas para compañías distribuidoras y surtidoras de Tegucigalpa. En el escenario analizado, el recorrido total pasó de 7,037.8 km a 4,008 km después de utilizar software de optimización y se estimó una reducción del 30 \% en los costos de distribución (p. 56). Más que tomar ese porcentaje como una meta para TravelDLS, el estudio sirve para mostrar que una propuesta de rutas puede evaluarse con datos concretos, como la distancia recorrida y los costos asociados.

\phantomsection\label{ant:hernandez}
Otro trabajo realizado en Honduras por \textcite{hernandez2024evaluacion} ejecutaron el proyecto de investigación titulado \textbf{Evaluación de alternativas de software de planificación de rutas para la empresa transporte Gonzales} en el cual se comparó distintas alternativas de software para la planificación de rutas en la empresa Transporte Gonzales. El estudio evaluó herramientas como Route4Me, Geotab y DispatchTrack y reportó una reducción calculada del 17 \% en los costos totales, menor que la meta inicial del 20 \% (pp. 10, 111). A la vez, la opción seleccionada permitía visualizar rutas, estado de envíos y tiempos estimados de entrega (p. 111). Este resultado es importante para TravelDLS porque recuerda que no conviene asumir de antemano que una herramienta será la mejor: la comparación debe hacerse con criterios claros y tomando en cuenta el contexto en el que se va a utilizar.

En Panamá, \textcite{castillo2025potenciando} se encontró una investigación titulada \textbf{Potenciando la competitividad logística de Panamá: una ingeniería inversa del transporte terrestre frente a Singapur y Rotterdam} analizó el transporte terrestre desde una perspectiva de competitividad logística e identificó, entre otros aspectos, costos elevados, limitaciones de infraestructura, baja digitalización y dificultades de coordinación (p. 10). Aunque el trabajo no desarrolla una aplicación informática, ayuda a entender que los problemas del transporte no dependen únicamente de la tecnología. TravelDLS puede apoyar la organización, el seguimiento y la planificación de los viajes, pero no pretende resolver por sí solo todos los factores que afectan al sector.

\phantomsection\label{ant:aragon}
En Nicaragua, \textcite{aragon2018implementacion} presentaron la investigación titulada \textbf{Implementación de un sistema de información para la administración de fletes terrestres en el departamento de transporte de la empresa master logistics, s.a}, cuyo propósito fue administrar fletes terrestres en dicha empresa. El trabajo buscó ordenar controles, mejorar la exactitud de los datos y facilitar el acceso a la información de las actividades logísticas; además, su objetivo general fue implementar un sistema capaz de registrar los elementos que forman parte de la administración de fletes (pp. 2, 9). Este antecedente es cercano a TravelDLS por su relación con el transporte de carga, aunque la propuesta actual amplía ese enfoque al incorporar contratación, seguimiento y planificación de rutas dentro de una plataforma pensada para distintas empresas.

\phantomsection\label{ant:garcia}
También en Nicaragua, \textcite{garcia2022desarrollo} presentaron la investigación titulada \textbf{Desarrollo de una aplicación móvil Android para el seguimiento y localización de las unidades de transporte público en la zona urbana del municipio de León, marzo - diciembre, año 2021} en la UNAN-León, la cual consistió en una aplicación Android para el seguimiento y localización de unidades de transporte público. El problema estudiado estuvo relacionado con la falta de información sobre recorridos, paradas y ubicación, y uno de sus objetivos fue mostrar en tiempo real la localización de las unidades mediante servicios de Google Maps (pp. 3, 7). Aunque se trata de transporte de pasajeros y no de carga pesada, el trabajo aporta una referencia nacional para el componente de ubicación de TravelDLS y demuestra que este tipo de tecnología ya se ha aplicado en proyectos de transporte desarrollados en el país.

Al observar los trabajos en conjunto, se aprecia que la administración del transporte, el seguimiento y la planificación de rutas han sido tratados desde diferentes enfoques. Los estudios internacionales muestran experiencias donde seguimiento y planificación forman parte de soluciones logísticas, mientras que los trabajos regionales permiten comparar resultados de optimización que no siempre son iguales. En Nicaragua se encontraron propuestas para administrar fletes y para localizar unidades de transporte, pero esas capacidades aparecen en proyectos separados. Con la revisión realizada hasta este momento, TravelDLS encuentra una oportunidad en reunir varias de estas funciones dentro de una plataforma orientada a la carga pesada. Esta idea todavía debe comprobarse con el mapeo sistemático completo, por lo que se mantiene como una brecha preliminar y no como una conclusión definitiva.

\subsection{Justificación}
\label{sec:justificacion}

TravelDLS se justifica porque las actividades de una operación de carga están relacionadas entre sí. Registrar un viaje tiene poco valor si luego esa información no puede conectarse con el vehículo, el conductor, la carga, la ruta y el estado del servicio. La propuesta busca que esos datos puedan consultarse de una manera más ordenada durante el recorrido y que cada actor tenga acceso a la información que le corresponde. La importancia de estas funciones también se observa en empresas logísticas que ya utilizan herramientas de seguimiento y planificación de rutas como parte de sus operaciones \parencite[p. 988]{pauzuoliene2024smart}.

El módulo de optimización se incorpora como un apoyo para planificar los recorridos, no como una promesa de que un algoritmo siempre encontrará la mejor opción para cualquier situación. Los antecedentes regionales muestran precisamente que los resultados cambian según los datos y la herramienta empleada: \textcite{aguilar2022optimizacion} obtuvo una reducción estimada del 30 \% en el escenario estudiado, mientras que \textcite{hernandez2024evaluacion} reportaron un 17 \% en su evaluación (p. 56; p. 111). Por eso, en TravelDLS será importante comparar los métodos definidos en el proyecto antes de valorar su comportamiento.

Sobre esta base, la innovación propuesta consiste en adaptar la relación entre contratación, gestión y seguimiento al transporte de carga pesada en Nicaragua. Frente a los antecedentes nacionales de administración de fletes para una empresa y localización de transporte de pasajeros \parencite{aragon2018implementacion,garcia2022desarrollo}, TravelDLS vinculará el servicio contratado con los recursos asignados por cada transportista y la consulta del estado de la carga por el cliente. La mejora esperada es dar continuidad a la información desde la contratación hasta el seguimiento del viaje. Como ya existen sistemas logísticos integrados y plataformas multitenant \parencite{sultan2022real,lasluisa2026diseno}, el aporte no se atribuye a esas tecnologías, sino a su aplicación al proceso descrito, cuya utilidad se evaluará en el proyecto. Los \textbf{\hyperref[sec:antecedentes-proyecto]{Antecedentes del Proyecto}} precisan las diferencias con cada estudio.

Para sostener esta propuesta, la plataforma contempla un modelo de pago en el que el porcentaje destinado al mantenimiento se agregará cuando la empresa transportista establezca el precio del viaje, de forma que el cliente vea desde el inicio el monto final antes de aceptar el servicio. Lo recaudado servirá para apoyar el funcionamiento de la plataforma, cubrir correcciones, mejoras y otros servicios necesarios para mantenerla disponible.

A partir de estas razones, resulta necesario delimitar de manera concreta qué componentes y funcionalidades formarán parte de TravelDLS durante esta etapa del proyecto. Por ello, a continuación se establece el alcance de la propuesta, definiendo hasta dónde llegará el desarrollo del sistema y qué elementos serán considerados dentro de su implementación.

\subsection{Alcance}
\label{sec:alcance}

TravelDLS se desarrollará como una plataforma web bajo el modelo SaaS, pensada para que distintas empresas de transporte de carga pesada puedan utilizarla. En esta primera etapa se contempla:

\begin{itemize}
	\item Registro y gestión de perfiles para empresas transportistas, conductores y clientes.
	\item Administración de vehículos, conductores, cargas, viajes, rutas y servicios asociados a cada empresa.
	\item Publicación de servicios de transporte y consulta o contratación por parte de los clientes.
	\item Asignación de conductores y vehículos a los viajes, con consulta del estado de cada operación.
	\item Localización y seguimiento de los viajes mediante servicios de mapas y funciones de ubicación disponibles en los dispositivos utilizados.
	\item Consulta para el cliente del estado y ubicación de su carga durante el desarrollo del servicio.
	\item Comparación de los resultados de Vecino Más Cercano y 2-opt con Google OR-Tools, siendo esta última la herramienta que se implementará en TravelDLS.
	\item Integración de una pasarela de pago en la que el porcentaje destinado al mantenimiento forme parte del precio final mostrado al cliente antes de aceptar el servicio.
	\item Realización de pruebas funcionales y elaboración de la documentación técnica y los manuales de usuario.
\end{itemize}

Aunque estas funcionalidades permiten establecer lo que se pretende desarrollar en TravelDLS, el proyecto también debe reconocer aquellos aspectos que no serán cubiertos durante esta primera etapa. Definir estos límites permite mantener el trabajo dentro de las condiciones académicas, técnicas y temporales disponibles.

\subsection{Limitaciones}
\label{sec:limitaciones}

Como TravelDLS se desarrollará en el marco de este proyecto académico y se concentrará en sus funciones principales, será necesario establecer algunos límites para mantener el trabajo enfocado y realizable:

\begin{itemize}
	\item La solución se desarrollará como aplicación web; en esta etapa no se contempla una aplicación móvil nativa independiente.
	\item El proyecto estará enfocado en el transporte terrestre de carga pesada y no abarcará transporte aéreo, marítimo, ferroviario ni transporte público de pasajeros.
	\item La optimización se probará en viajes con varios destinos y con los métodos definidos para el proyecto, por lo que no se pretende cubrir todos los tipos de problemas de rutas que existen.
	\item La actualización de la ubicación dependerá de la conexión a internet, del dispositivo utilizado y de los servicios de ubicación disponibles.
	\item Las pruebas se realizarán con los escenarios disponibles dentro del proyecto académico, por lo que sus resultados no representan todavía una evaluación de uso a gran escala.
	\item En esta etapa no se incluirán conexiones especializadas con sistemas externos como plataformas contables, aduaneras, aseguradoras u otros servicios que requieran desarrollos adicionales.
	\item TravelDLS se plantea para ser utilizado por distintas empresas transportistas y no como una solución exclusiva para una sola organización.
\end{itemize}

\subsection{Viabilidad del proyecto}
\label{sec:viabilidad}

Una vez definido lo que formará parte del proyecto y los límites bajo los cuales será desarrollado, corresponde analizar si la propuesta puede llevarse a la práctica con los recursos, tecnologías y condiciones disponibles. Para ello, la viabilidad de TravelDLS se estudia desde diferentes perspectivas.

\subsubsection{Viabilidad Técnica}

TravelDLS presenta viabilidad técnica debido a que las tecnologías seleccionadas permiten desarrollar una plataforma web con una separación clara entre la interfaz de usuario, la lógica del servidor, el almacenamiento de la información, la gestión de archivos multimedia y la visualización geográfica generada por empresas transportistas, conductores y clientes. Para el frontend se utilizará Next.js, mientras que el backend será desarrollado con AdonisJS. La información será almacenada en PostgreSQL, las fotografías de usuarios y empresas y el logotipo de la aplicación se gestionarán mediante Cloudinary, y los mapas, rutas y seguimiento de ubicación se apoyarán en Mapbox. Además, se utilizarán JSON Web Tokens (JWT) para apoyar los procesos de autenticación y control de acceso. El frontend será desplegado en Vercel y el backend en Render, utilizando un dominio adquirido mediante GoDaddy o Namecheap y conexiones HTTPS protegidas mediante certificados SSL/TLS.

Como referencia para la organización interna del backend, se tomó en cuenta el trabajo de \textcite{arce2026arquitectura}, quienes desarrollaron una plataforma web utilizando una estructura basada en Modelo-Vista-Controlador (MVC) complementada con patrones como Service Layer, con el propósito de separar responsabilidades entre los distintos componentes del sistema y facilitar su mantenimiento (pp.~28, 33--35). Este antecedente resulta útil para TravelDLS porque muestra una forma de mantener la lógica de negocio separada de los controladores. Sin embargo, la arquitectura de TravelDLS no busca reproducir completamente la estructura utilizada en dicho trabajo, sino aplicar específicamente MVC + Service Layer dentro del backend desarrollado con AdonisJS.

En el backend de TravelDLS se implementará una estructura basada en el patrón Modelo-Vista-Controlador (MVC), complementada con una Service Layer. Esta organización se aplicará internamente en AdonisJS y no implica que Next.js forme parte de la Vista del MVC del backend. El Modelo representará las entidades y relaciones persistidas en PostgreSQL, mientras que los Controladores recibirán las solicitudes provenientes del frontend y coordinarán las respuestas de la API. La lógica principal del negocio se trasladará a la capa de servicios, donde se concentrarán procesos como la gestión de viajes, la asignación de vehículos y conductores, la contratación de servicios, el seguimiento y la optimización de rutas. En una plataforma web estudiada por \textcite{arce2026arquitectura}, la combinación de MVC con patrones como Service Layer permitió separar responsabilidades y favorecer la mantenibilidad y escalabilidad del sistema. A partir de este enfoque, TravelDLS utilizará MVC + Service Layer en el backend para mantener los controladores ligeros y separar las reglas del negocio de la gestión de solicitudes.

Next.js se utilizará como aplicación frontend independiente encargada de la interfaz de usuario y de consumir la API expuesta por AdonisJS. Esta separación permitirá desarrollar vistas específicas para empresas transportistas, conductores y clientes sin mezclar la lógica de presentación con las reglas del negocio. AdonisJS, por su parte, concentrará la lógica del servidor y facilitará la organización de módulos relacionados con usuarios, empresas, vehículos, cargas, viajes, rutas, servicios y pagos.

PostgreSQL se selecciona porque la información de TravelDLS posee una naturaleza principalmente relacional. Las empresas estarán asociadas con usuarios, vehículos, conductores, cargas y viajes, mientras que cada viaje se relacionará con rutas, estados y clientes. Por ello, una base de datos relacional resulta adecuada para conservar estas dependencias y mantener consistencia entre los registros. JWT se utilizará como mecanismo de apoyo para la autenticación entre el frontend y el backend, complementado con controles de autorización que permitan diferenciar los permisos de cada tipo de usuario.

Las fotografías de los usuarios y de las empresas, así como el logotipo de TravelDLS, se alojarán en Cloudinary para facilitar su carga, organización y visualización dentro de la plataforma. En PostgreSQL se conservarán las referencias necesarias para asociar cada imagen con su usuario, empresa o elemento de la aplicación. Esta integración tiene antecedentes en TaxiGo, que utiliza Cloudinary para gestionar archivos multimedia \parencite{yaqoub2024taxigo}, y en Truxi, donde las imágenes se almacenan en ese servicio y sus enlaces se relacionan con los registros del sistema \parencite{vad2020truxi}. Para los mapas se utilizará Mapbox, que permitirá representar las rutas y la ubicación de los vehículos durante los viajes. Su aplicación en soluciones de transporte y logística se observa en el sistema de seguimiento en tiempo real desarrollado por \textcite{aronsson2021tracking}, donde Mapbox se seleccionó como proveedor de mapas, y en Thunder Logistics, que integra Mapbox GL JS, Directions API y Geocoding API para mapas interactivos, rutas y seguimiento público \parencite{deyheralde2026thunder}. La posición se obtendrá de los servicios de ubicación del dispositivo del conductor, se enviará al backend y se actualizará en el mapa conforme se reciban nuevos datos; de esta manera, el seguimiento podrá mostrar la ubicación más reciente disponible.

Para el despliegue, Vercel alojará el frontend desarrollado en Next.js y Render alojará el backend construido con AdonisJS. Esta distribución permitirá actualizar ambos componentes de forma independiente y mantener la aplicación disponible desde una infraestructura centralizada. El dominio será adquirido mediante GoDaddy o Namecheap y se utilizará HTTPS mediante certificados SSL/TLS para proteger la comunicación entre los usuarios y los servicios del sistema.

En conjunto, la propuesta se considera técnicamente viable porque las tecnologías seleccionadas permiten separar la interfaz, la lógica del servidor y la persistencia de datos de acuerdo con las necesidades definidas para TravelDLS. La adopción de MVC + Service Layer específicamente en el backend permitirá ordenar el código y separar responsabilidades, mientras que la división entre frontend, backend y base de datos facilitará el mantenimiento y futuras ampliaciones de la plataforma.

Contar con tecnologías adecuadas y una estructura definida permite establecer que el sistema puede construirse desde el punto de vista técnico. Sin embargo, también es necesario analizar cómo estas herramientas podrán utilizarse dentro de las actividades diarias de empresas, conductores y clientes, lo que conduce al estudio de la viabilidad operativa.

\subsubsection{Viabilidad Operativa}

La viabilidad operativa de TravelDLS se basa en que será una plataforma web accesible mediante un navegador, por lo que las empresas transportistas no necesitarán instalar una versión independiente del sistema ni mantener servidores propios. El acceso estará dividido principalmente entre empresas transportistas, conductores y clientes, mostrando a cada usuario las funciones relacionadas con su participación dentro del servicio.

La utilización de sistemas digitales con diferentes tipos de usuarios cuenta con antecedentes dentro del área logística. \textcite{sultan2022real} desarrollaron una solución móvil y web que integra clientes, conductores y administradores mediante funciones de registro, reservas, seguimiento y gestión logística (pp.~1--4). Este antecedente es relevante para TravelDLS porque demuestra que los procesos de transporte pueden organizarse mediante una plataforma donde varios actores interactúan con información relacionada con un mismo servicio.

La utilidad operativa del seguimiento y la planificación también ha sido documentada en empresas logísticas. \textcite{pauzuoliene2024smart} encontraron que, dentro de la muestra analizada, el 69.6\,\% de las empresas utilizaba tecnologías para seguimiento de carga en tiempo real y el 60.9\,\% empleaba herramientas para planificación o selección de rutas (pp.~988--989). Aunque estos resultados corresponden a un contexto distinto al nicaragüense, muestran que ambas funciones forman parte de procesos tecnológicos utilizados en el sector logístico.

En TravelDLS, las empresas podrán administrar vehículos, conductores, cargas, viajes y servicios; los conductores podrán consultar la información de los recorridos asignados; y los clientes podrán conocer el estado de los servicios contratados. Al mantenerse la aplicación centralizada, las actualizaciones realizadas sobre la plataforma estarán disponibles para todos los usuarios sin requerir instalaciones adicionales en cada empresa.

Como parte de esa gestión, los usuarios y las empresas podrán incorporar sus fotografías a los perfiles y registros que les correspondan, mientras que el logotipo de TravelDLS se mostrará de forma consistente en la interfaz. La experiencia de TaxiGo muestra la utilidad de integrar la gestión de imágenes y las funciones de seguimiento dentro de un sistema de transporte con distintos tipos de usuarios \parencite{yaqoub2024taxigo}. Durante un viaje activo, el conductor compartirá su ubicación mediante el dispositivo que utilice y Mapbox la presentará sobre el mapa para que la empresa y el cliente consulten el avance del servicio según sus permisos. Truxi ofrece un antecedente de uso de mapas para indicar lugares de recogida y entrega y apoyar los recorridos de los conductores \parencite{vad2020truxi}; además, Chalmers documenta el uso de Mapbox dentro de un sistema de seguimiento posicional en tiempo real para una empresa de transporte \parencite{aronsson2021tracking}, y Thunder Logistics lo emplea para visualizar rutas, envíos activos y páginas de tracking con ubicación actualizada \parencite{deyheralde2026thunder}. En TravelDLS, esta información se vinculará con el viaje registrado, de modo que el seguimiento y la consulta de la ruta formen parte del trabajo cotidiano de los participantes.

El funcionamiento de TravelDLS dependerá de factores como la conexión a Internet, la disponibilidad de los servicios de ubicación y los dispositivos utilizados por los conductores. Por ello, estos elementos se considerarán durante el diseño de la plataforma, tomando como referencia investigaciones previas y documentación técnica relacionada con la digitalización de los procesos logísticos. Asimismo, la pertinencia operativa de la solución se analizará mediante los procedimientos de diseño y evaluación establecidos en la metodología DSRM, con el propósito de determinar si las funcionalidades propuestas responden al problema identificado y a los requerimientos definidos para el sistema.

Además de comprobar que TravelDLS puede integrarse a las actividades de sus posibles usuarios, es necesario determinar si su desarrollo y funcionamiento pueden sostenerse económicamente. Por esta razón, el siguiente análisis considera la inversión inicial, los costos operativos y los ingresos que podrían generarse mediante el modelo propuesto para la plataforma.

\subsubsection{Viabilidad Económica}

La inversión inicial de TravelDLS se calcula sumando el desarrollo del sistema, el dominio del primer año y la configuración, las pruebas y la puesta en marcha: C\$52,732.80 + C\$439.44 + C\$3,662.00 = C\$56,834.24. Para expresar en córdobas los importes presupuestados en dólares, se multiplica cada monto por C\$36.62. Este valor corresponde al redondeo a dos decimales del tipo de cambio oficial de C\$36.6243 por USD~1.00 publicado por el Banco Central de Nicaragua para 2026 \parencite{bcn2026tipocambio}. A continuación se explica la procedencia de los rubros que integran esta estimación académica.

Para el desarrollo del sistema se asigna una valoración académica de C\$17,577.60 por integrante. Al multiplicar este monto por los tres integrantes se obtiene C\$52,732.80. Como referencias salariales, World Salaries reporta un promedio mensual de C\$18,158 para desarrolladores y programadores en Nicaragua, mientras que WageIndicator presenta un rango mensual inicial de C\$15,612 a C\$31,324 \parencite{worldsalaries2026developer,wageindicator2026desarrolladores}. Estas referencias orientan la valoración del trabajo técnico; el monto adoptado por integrante es una estimación académica por dedicación parcial al proyecto. Se contabiliza como parte de la inversión por el valor del trabajo aportado, aunque no se desembolse como salario.

Para el dominio del primer año se presupuestan USD~12.00, equivalentes a C\$439.44 al multiplicar $12.00 \times 36.62$. Las referencias de Namecheap y GoDaddy permiten consultar precios de registro y renovación, que varían según el dominio, las promociones y el período contratado \parencite{namecheapdomain,godaddy2026domain}. Por ello, los USD~12.00 se utilizan como una asignación presupuestaria aproximada y no como una cotización definitiva. La configuración, las pruebas y la puesta en marcha se valoran internamente en USD~100.00, equivalentes a C\$3,662.00 mediante $100.00 \times 36.62$; este rubro corresponde a una estimación del equipo para preparar el funcionamiento de la plataforma.

GitHub Free y los certificados SSL/TLS de Let's Encrypt se incluyen con un costo inicial de C\$0.00, de acuerdo con sus servicios gratuitos publicados \parencite{githubpricing,letsencrypt2021about}. Asimismo, el presupuesto inicial contempla Cloudinary para las imágenes y Mapbox para los mapas dentro de sus niveles gratuitos \parencite{cloudinarypricing,mapboxpricing}. El costo de C\$0.00 se mantiene mientras el uso permanezca dentro de los límites aplicables de cada servicio. La Tabla~\ref{tab:inversion-inicial} reúne los rubros de la inversión inicial.

\begin{table}[H]
\centering
\caption{Inversión inicial estimada de TravelDLS}
\label{tab:inversion-inicial}
\begin{tabularx}{\textwidth}{>{\RaggedRight\arraybackslash}X r}
\toprule
Concepto & Monto (C\$) \\
\midrule
Desarrollo del sistema (tres integrantes) & 52,732.80 \\
Dominio del primer año & 439.44 \\
Configuración, pruebas y puesta en marcha & 3,662.00 \\
GitHub Free y certificado SSL/TLS de Let's Encrypt & 0.00 \\
Cloudinary y Mapbox dentro de sus niveles gratuitos & 0.00 \\
\midrule
Total & 56,834.24 \\
\bottomrule
\end{tabularx}
\end{table}

Los costos operativos se calculan a partir del alojamiento de la aplicación, la base de datos, el soporte y una reserva para consumos adicionales. Para el backend se toma como referencia una instancia de Render de USD~7.00 mensuales, equivalentes a C\$256.34 mediante $7.00 \times 36.62$ \parencite{renderpricing}. Para PostgreSQL se conserva una asignación presupuestaria de USD~10.00 mensuales, equivalente a C\$366.20; su costo definitivo deberá confirmarse según la capacidad y el almacenamiento contratados en Render. Para el frontend se considera la tarifa base de Vercel Pro de USD~20.00 mensuales, equivalente a C\$732.40 mediante $20.00 \times 36.62$, bajo el supuesto de una sola plaza de desarrollador de pago \parencite{vercelpricing}. El soporte y mantenimiento se estiman internamente en USD~100.00 mensuales, equivalentes a C\$3,662.00, para atender ajustes y tareas de seguimiento técnico.

Aunque se prevé iniciar con los niveles gratuitos de Mapbox y Cloudinary, el uso de mapas, solicitudes, almacenamiento y procesamiento de imágenes puede generar cargos al superar los límites de cada servicio \parencite{mapboxpricing,cloudinarypricing}. Para contemplar esa posibilidad, la Tabla~\ref{tab:costos-operativos} incluye una \textit{Reserva para consumos adicionales} de C\$732.40 mensuales, equivalente a USD~20.00. Su valor anual se obtiene multiplicando C\$732.40 por 12 meses, con un resultado de C\$8,788.80. Esta reserva ya forma parte del total de costos operativos y está destinada a atender posibles cargos adicionales o contribuir a la contratación de un plan de pago. Su suficiencia dependerá del consumo real y de las tarifas aplicables; si el costo supera lo reservado, deberá ajustarse el presupuesto.

La suma de los rubros mensuales es C\$256.34 + C\$366.20 + C\$732.40 + C\$3,662.00 + C\$732.40 = C\$5,749.34. Al multiplicar este total por 12 se obtienen C\$68,992.08 anuales. La Tabla~\ref{tab:costos-operativos} presenta el desglose mensual y anual, incluida la reserva.

\begin{table}[H]
\centering
\caption{Costos operativos estimados de TravelDLS}
\label{tab:costos-operativos}
\begin{tabularx}{\textwidth}{>{\RaggedRight\arraybackslash}X r r}
\toprule
Concepto & Mensual (C\$) & Anual (C\$) \\
\midrule
Backend en Render & 256.34 & 3,076.08 \\
PostgreSQL en Render & 366.20 & 4,394.40 \\
Frontend en Vercel Pro & 732.40 & 8,788.80 \\
Soporte y mantenimiento & 3,662.00 & 43,944.00 \\
Reserva para consumos adicionales & 732.40 & 8,788.80 \\
\midrule
Total & 5,749.34 & 68,992.08 \\
\bottomrule
\end{tabularx}
\end{table}

El dominio del primer año se incluye en la inversión inicial para no duplicarlo como gasto operativo. A partir del segundo año, el modelo agrega una renovación anual estimada de C\$439.44, por lo que los costos anuales proyectados ascienden a C\$69,431.52. Esta renovación conserva el supuesto de USD~12.00 del presupuesto, aunque las tarifas de renovación pueden diferir de las de registro y deberán confirmarse con el proveedor \parencite{namecheapdomain,godaddy2026domain}. La Tabla~\ref{tab:supuestos-economicos} resume los valores utilizados como base del modelo.

\begin{table}[H]
\centering
\caption{Supuestos base del modelo económico}
\label{tab:supuestos-economicos}
\begin{tabularx}{\textwidth}{>{\RaggedRight\arraybackslash}X r}
\toprule
Concepto & Valor (C\$) \\
\midrule
Inversión inicial & 56,834.24 \\
Costos operativos anuales & 68,992.08 \\
Renovación de dominio anual & 439.44 \\
\bottomrule
\end{tabularx}
\end{table}

Con estos costos como base, la viabilidad económica se proyecta a cinco años. El modelo de ingresos considera un recargo del 8\,\% sobre el precio definido por la empresa transportista. Para un envío base de C\$3,662.00, el cliente pagaría C\$3,954.96; de ese monto, C\$292.96 corresponden al ingreso bruto de TravelDLS. Al descontar una comisión estimada del 3\,\% por pasarela de pago, equivalente a C\$118.65, queda un ingreso neto de C\$174.31 por envío. Estos porcentajes son supuestos del modelo académico. El detalle del cálculo por envío se presenta en la Tabla~\ref{tab:anexo-contribucion-envio} del Anexo~\ref{anexo:flujo-caja}.

La proyección considera 60 envíos mensuales en el primer año y un aumento de 30 envíos mensuales por cada año posterior, hasta llegar a 180 envíos mensuales en el quinto año. El flujo completo se presenta en la Tabla~\ref{tab:anexo-flujo-caja-proyectado} del Anexo~\ref{anexo:flujo-caja}.

Con una tasa académica de descuento del 12\,\%, el modelo obtiene un valor actual neto (VAN) de C\$547,106.97, una tasa interna de retorno (TIR) de 162\,\% y un período de recuperación de la inversión (PRI) de 1.002 años. Estos indicadores se muestran en la Tabla~\ref{tab:anexo-indicadores-financieros} del Anexo~\ref{anexo:flujo-caja}. Bajo estos supuestos, TravelDLS presenta viabilidad económica preliminar, aunque los resultados deberán revisarse cuando se conozcan la comisión definitiva de pago, la demanda real, la renovación del dominio y el consumo efectivo de infraestructura.

No obstante, la sostenibilidad económica no constituye el único aspecto que debe considerarse antes de implementar TravelDLS. Debido a que la plataforma gestionará datos personales, ubicaciones, contrataciones y servicios relacionados con el transporte terrestre, también debe revisarse el marco legal aplicable a su funcionamiento.

\subsubsection{Viabilidad Legal}

La viabilidad legal de TravelDLS se sustenta en las normas nicaragüenses que regulan la protección de datos, el comercio electrónico y el transporte terrestre de carga. La Ley No. 787, Ley de Protección de Datos Personales, es necesaria porque la plataforma manejará información de empresas, conductores y clientes, incluyendo ubicaciones en tiempo real, fotografías de perfil y documentos de identidad. Su artículo 4 condiciona la creación de ficheros de datos personales a que exista autorización y consentimiento del titular, lo cual obliga a TravelDLS a implementar mecanismos de consentimiento informado durante el registro de usuarios. El artículo 5 pide que los datos sean los necesarios para el fin que se persigue, de modo que la plataforma no puede recolectar información más allá de la que realmente utiliza para gestionar los viajes. El artículo 7 obliga a informar al titular para qué se usarán sus datos, quién los verá y cómo puede ejercer sus derechos de acceso, rectificación o cancelación \parencite[Arts.4,5 y 7]{ley787_2012}. Sin esta ley, TravelDLS no tendría un marco claro para tratar la información de los usuarios, y las empresas transportistas podrían exponerse a reclamos por el mal uso de datos de sus conductores o clientes.

La contratación de servicios a través de TravelDLS también tiene respaldo legal. La Ley No. 729, Ley de Firma Electrónica, es necesaria porque la plataforma está diseñada para que empresas transportistas publiquen servicios y los clientes los contraten directamente desde la web, sin necesidad de firmar documentos en físico. Su artículo 1 otorga eficacia y valor jurídico a la firma electrónica y a los certificados digitales, así como a toda información en formato electrónico. El artículo 6 establece que la firma electrónica certificada tendrá el mismo valor jurídico que la firma manuscrita y será admisible como medio de prueba en procesos judiciales o administrativos \parencite[Arts. 1 y 6]{ley729_2010}. Esto significa que los contratos celebrados a través de la plataforma entre empresas, conductores y clientes tendrán pleno valor jurídico, y permiten que la contratación dentro de TravelDLS no dependa de trámites presenciales. La Ley 729 fue publicada en La Gaceta No. 165 del 30 de agosto de 2010, y es la norma vigente que regula esta materia en Nicaragua.

En cuanto a la seguridad de la información, el Código Penal, Ley No. 641, es necesario porque TravelDLS operará como una plataforma en línea que almacena información crítica para la operación de las empresas transportistas. Su artículo 246 castiga con prisión de uno a tres años a quien, con intención de dañar, distribuya programas informáticos destructivos, y su artículo 245 sanciona la destrucción o inutilización de registros informáticos, con penas mayores si se trata de información ligada a un servicio público \parencite[Arts.245--246]{ley641_2008}. Esto obliga a TravelDLS a proteger su infraestructura con medidas como los tokens JWT y las conexiones HTTPS que ya se mencionan en la propuesta técnica, y justifica la inversión en seguridad para evitar que un ataque informático deje a las empresas sin acceso a sus viajes, rutas o datos de clientes.

En el ámbito del transporte, la Ley No. 524, Ley General de Transporte Terrestre, y su Reglamento (Decreto No. 42-2005), son necesarios porque TravelDLS está dirigido específicamente al transporte terrestre de carga pesada. El artículo 5 de la Ley 524 clasifica el transporte de carga como el traslado de bienes y todo tipo de objetos, sólidos o líquidos, que se presta a otros. El artículo 7 exige que todos los vehículos de servicio público de transporte terrestre cumplan las normas técnicas de inspección y obtengan la constancia de aprobación correspondiente \parencite[Arts. 5 y 7]{ley524_2005}. El Reglamento, en su artículo 13, detalla que la inspección técnica mecánica es el proceso mediante el cual una unidad de transporte público se somete a revisión de funcionamiento y estado de sus partes, y debe realizarse por lo menos cada seis meses \parencite[Art.13]{decreto42_2005}. Aunque TravelDLS no es una empresa transportista, sus módulos de registro de vehículos, conductores y viajes deben alinearse con estos requisitos para que la plataforma realmente sirva de apoyo a las empresas en el cumplimiento de la normativa, y no se convierta en un sistema que ignore las obligaciones legales del sector.

La Ley No. 842, Ley de Protección de los Derechos de las Personas Consumidoras y Usuarias, también tiene relevancia. Es necesaria porque los clientes que contraten transporte de carga a través de TravelDLS podrían quedar en una situación de desventaja informativa frente a la empresa transportista. Su artículo 2 cubre todos los actos de consumo dentro del país e incluso protege a microempresas cuando hay asimetría informativa con el proveedor. El artículo 3 excluye a quienes compran para su actividad empresarial, pero los clientes que contraten servicios de carga podrían no tener el mismo conocimiento técnico que la empresa sobre rutas, tiempos o costos \parencite[Arts. 2--3]{ley842_2013}. Por eso, la plataforma debe mostrar de forma clara el precio final, incluyendo el recargo del 8 \% y las condiciones del servicio, para que el cliente sepa exactamente qué está contratando antes de aceptar.

Estas leyes no solo dan respaldo jurídico a TravelDLS, sino que definen cómo debe operar la plataforma en aspectos concretos como el consentimiento de datos, la validez de los contratos electrónicos, la seguridad de la información, el cumplimiento normativo del transporte y la protección del cliente. Sin este marco, la plataforma podría funcionar técnicamente, pero quedaría expuesta a reclamos legales y no sería una herramienta confiable para las empresas transportistas que la utilicen.

Una vez consideradas las condiciones técnicas, operativas, económicas y legales que rodean la propuesta, también resulta necesario observar las condiciones en las que los posibles usuarios interactuarán con la plataforma. El funcionamiento de TravelDLS no dependerá únicamente del cumplimiento de requisitos técnicos o normativos, sino también de factores relacionados con la adopción de tecnología, las formas habituales de trabajo, la comunicación entre los participantes, las condiciones de acceso digital y la confianza en la información proporcionada. Por esta razón, a continuación se analizan algunos aspectos socioculturales y tecnológicos que deben tomarse en cuenta durante el diseño de la plataforma.

\subsection{Aspectos socioculturales}
\label{sec:aspectos-socioculturales}

El desarrollo de TravelDLS no depende únicamente de la implementación de sus componentes tecnológicos, sino también de las condiciones bajo las cuales las empresas transportistas, conductores y clientes podrían utilizar la plataforma. En las MIPYMES latinoamericanas, la transformación digital presenta diferentes niveles de avance según el tamaño de la empresa, los recursos disponibles y su grado de madurez tecnológica \parencite{dini2021transformacion}. Por esta razón, aspectos como la adopción tecnológica, los hábitos de trabajo, la comunicación, el acceso digital y la confianza deben considerarse durante el diseño de una solución dirigida al transporte terrestre de carga pesada.

\subsubsection{Adopción de la tecnología}
\label{sec:adopcion-tecnologia}

La incorporación de nuevas herramientas digitales puede representar un cambio importante para las MIPYMES, especialmente cuando modifica procedimientos que los trabajadores ya realizan de una manera conocida. \textcite{dini2021transformacion} señalan que la digitalización de las pequeñas empresas no ocurre de forma homogénea, debido a diferencias en sus capacidades, recursos y nivel de desarrollo tecnológico. En Nicaragua, \textcite{barreda2023economia} encontraron que MIPYMES de Estelí ya utilizan dispositivos y herramientas digitales como computadoras, teléfonos inteligentes y redes sociales dentro de sus actividades, lo que muestra una incorporación progresiva de la tecnología en pequeñas empresas del país.

Ante este contexto, TravelDLS busca facilitar su adopción mediante una plataforma web con procesos organizados y una interfaz comprensible para sus distintos usuarios. Las funciones de registro de vehículos, conductores, cargas, viajes y rutas deben mantener una estructura consistente, evitando que la digitalización introduzca una complejidad mayor que los procedimientos que pretende mejorar. De esta manera, la usabilidad y una curva de aprendizaje razonable se consideran elementos importantes para favorecer la integración progresiva de la plataforma en las actividades de las empresas transportistas.

\subsubsection{Hábitos de trabajo}
\label{sec:habitos-trabajo}

La digitalización de los procesos logísticos no consiste únicamente en sustituir registros físicos por formularios electrónicos, sino en mejorar la forma en que la información se organiza y se relaciona durante una operación. \textcite{aleman2021sistema}, en un estudio desarrollado para la compañía Logística de Carga Intermodal en Managua, identificaron que la empresa utilizaba diferentes herramientas computacionales para procesos como cotizaciones, inventario y facturación, pero estas no satisfacían completamente sus necesidades, lo que motivó el desarrollo de un sistema integrado de gestión.

Esta situación guarda relación con TravelDLS, cuyo propósito es evitar que la información de una operación se gestione como elementos independientes. La plataforma busca relacionar la empresa, el cliente, la carga, el vehículo, el conductor, el viaje y la ruta dentro de un mismo servicio. Ante esta necesidad, TravelDLS responde mediante la centralización de información y la reutilización de datos previamente registrados, permitiendo que las actividades digitales sigan el flujo de trabajo de la operación y reduzcan la necesidad de mantener registros separados.

\subsubsection{Comunicación entre los actores}
\label{sec:comunicacion-actores}

Las operaciones de transporte requieren coordinación entre participantes que normalmente se encuentran en lugares diferentes. La empresa administra el servicio, el conductor realiza el recorrido y el cliente necesita conocer el avance de su carga. \textcite{heinbach2022datadriven} muestran que las plataformas digitales aplicadas al transporte terrestre pueden integrar información relacionada con órdenes, recursos y seguimiento, facilitando la visibilidad de las operaciones y el intercambio de información entre los participantes del proceso.

Ante esta realidad, TravelDLS busca disminuir la dependencia de comunicaciones individuales para consultar información básica del servicio. La empresa podrá conocer el estado de sus operaciones y recursos, el conductor tendrá acceso a la información correspondiente al viaje asignado y el cliente podrá consultar el estado y ubicación de su carga. Esto no pretende reemplazar completamente la comunicación directa, sino disponer de una fuente central de información que permita a cada actor consultar los datos que le corresponden según su rol.

\subsubsection{Acceso digital}
\label{sec:acceso-digital}

La disponibilidad de conectividad constituye un factor crítico a considerar, especialmente en una plataforma que integra funciones de localización y seguimiento. Según datos de la Unión Internacional de Telecomunicaciones, en Nicaragua existe una cobertura significativa de redes móviles; para el año 2022, se reportaba una cobertura de al menos 3G del 85~\% y de 4G del 73~\%, aunque con disparidades entre tecnologías \parencite{ituDigitalNicaragua}. No obstante, la existencia de cobertura no garantiza una calidad de servicio constante. TELCOR reportó para 2024 una cobertura móvil nacional del 87,7~\%, señalando que varios indicadores de calidad requieren mejoras \parencite{telcor2024calidad}. Ante este panorama, TravelDLS debe prever que la conexión puede variar a lo largo de un recorrido. Por ello, la plataforma se concibe como una aplicación web adaptable a diferentes tamaños de pantalla, diseñada para gestionar adecuadamente los retrasos o interrupciones en las actualizaciones de ubicación. Asimismo, la visualización de la fecha y hora de la última posición recibida resulta fundamental para evitar que una ausencia temporal de datos sea interpretada erróneamente como una detención del vehículo.

\subsubsection{Transparencia y confianza}
\label{sec:transparencia-confianza}

La confianza en un servicio de transporte también se relaciona con la calidad y disponibilidad de la información durante su ejecución. \textcite{nie2026visibility} distinguen entre visibilidad, trazabilidad y transparencia: la visibilidad se relaciona con disponer de información actual sobre una operación, la trazabilidad con conservar información que permita reconstruir lo ocurrido y la transparencia con proporcionar información adecuada a los actores involucrados.

Esta relación entre información y confianza también se observa en el estudio de \textcite{brookbanks2022impact}, quienes encontraron que una plataforma digital basada en datos compartidos y confiables puede reducir la duplicación de información, mejorar la visibilidad de la cadena de suministro y apoyar la construcción de confianza entre sus participantes. Los autores también aclaran que la tecnología no elimina la necesidad de establecer relaciones de confianza entre los actores involucrados.

En TravelDLS, estos principios se reflejan en la consulta del estado y ubicación de los servicios. La posición más reciente del vehículo aporta visibilidad, mientras que el registro de estados y ubicaciones contribuye a la trazabilidad de la operación. A su vez, el cliente podrá consultar la información correspondiente al servicio contratado, mientras que los mecanismos de autenticación y roles limitarán el acceso según las responsabilidades de cada usuario. De esta manera, la transparencia no implica mostrar toda la información disponible, sino proporcionar datos relevantes y confiables sin comprometer la privacidad de los participantes.

En conjunto, estos aspectos muestran que la incorporación de TravelDLS no debe analizarse únicamente desde su funcionamiento técnico, sino también desde la manera en que sus posibles usuarios podrían relacionarse con la plataforma durante sus actividades cotidianas. La adopción progresiva de herramientas digitales, la organización de los procesos de trabajo, la coordinación entre empresa, conductor y cliente, las condiciones variables de conectividad y la necesidad de disponer de información confiable influyen directamente en las decisiones de diseño del sistema. Por esta razón, funcionalidades como la centralización de la información, la organización por roles, una interfaz adaptable y el seguimiento de las operaciones no responden solamente a requerimientos funcionales, sino también a las condiciones identificadas en la literatura sobre el uso de tecnologías digitales. Considerados estos factores, resulta necesario profundizar en los conceptos y fundamentos que explican los componentes principales de TravelDLS, los cuales se desarrollan a continuación en el Marco Teórico.